% Riemann--Siegel recursion / Hardy Z(t) zero-transfer certificates.
\subsubsection{Riemann--Siegel 余项的层级递归与零点漂移证书：Hardy 实值化的局部传输}\label{subsubsec:xi-riemann-siegel-recursive-zero-transfer}
上一小节以 Poisson 混叠把离散采样误差转写为别名尾和，并由 Rouch\'e 定理将“函数误差预算”转换为“零点漂移预算”（定理 \ref{thm:xi-dyadic-zero-tracking}）。
本小节记录另一条与之同型的局部证书链：在 Hardy 实值化与 Riemann--Siegel 余项的高阶渐近展开下，截断阶数 \(m\) 形成一条自然的递归层级；通过一个 Rouch\'e--Taylor 传输引理，可把余项预算直接译码为零点漂移界与区间内零点计数的闭合条件。

\paragraph{Hardy 实值化与递归截断}
记 Riemann--Siegel \(\vartheta\)-函数为 \(\vartheta(t)\)，并定义 Hardy 实值化函数
\begin{equation}\label{eq:xi-hardy-Z-def}
Z(t):=e^{\mathrm{i}\vartheta(t)}\zeta\!\left(\tfrac12+\mathrm{i}t\right)\qquad (t\in\RR).
\end{equation}
则 \(Z(t)\in\RR\) 且其实零点与 \(\zeta(\tfrac12+\mathrm{i}t)\) 的临界线零点一致。
为避免与前文的双变量配分函数 \(Z_m(y)\) 混淆，本段的 \(Z(t)\) 专指 \eqref{eq:xi-hardy-Z-def} 所定义的 Hardy 实值化函数。
令
\begin{equation}\label{eq:xi-riemann-siegel-aNz}
a(t):=\sqrt{\frac{t}{2\pi}},\qquad N:=\lfloor a(t)\rfloor,\qquad z:=1-2\bigl(a(t)-N\bigr),
\end{equation}
并以 \(R(t)\) 记 Riemann--Siegel 公式中的余项。
在标准高阶渐近展开口径中，可将余项写为
\begin{equation}\label{eq:xi-riemann-siegel-remainder-expansion}
R(t)=(-1)^{N+1}a(t)^{-1/2}\sum_{r\ge 0} C_r\!\bigl(z(t)\bigr)\,a(t)^{-r},
\end{equation}
其中 \(C_r\) 为一族由局部参数 \(z\) 决定的显式系数函数。
令 \(Z^{(m)}(t)\) 表示把 \eqref{eq:xi-riemann-siegel-remainder-expansion} 截断到 \(r=0,\dots,m-1\) 所得到的递归截断近似（其余 Riemann--Siegel 主项取同一约定）。
在固定截断索引与分支的局部邻域内，\(t\mapsto Z^{(m)}(t)\) 作为复变量具有全纯延拓；在下述论证中，仅使用该局部全纯性与一致误差预算，而不依赖余项系数的具体构造。

\paragraph{Rouch\'e--Taylor 零点传输引理}
\begin{theorem}[Rouch\'e--Taylor 递归零点传输：从可控余项到可控零点漂移]\label{thm:xi-rouche-taylor-zero-transfer}
设 \(f,g\) 在闭圆盘 \(\overline{D(t_0,r)}:=\{t\in\CC:\ \abs{t-t_0}\le r\}\) 上全纯，且 \(g(t_0)=0\)、\(g'(t_0)\neq 0\)。
记
\[
\varepsilon:=\sup_{\abs{t-t_0}\le r}\abs{f(t)-g(t)},
\qquad
M:=\sup_{\abs{t-t_0}\le r}\abs{g''(t)}.
\]
若满足
\begin{equation}\label{eq:xi-rouche-taylor-conditions}
0<r,\quad
\abs{g'(t_0)}>\tfrac{M}{2}r\quad\text{且}\quad
\varepsilon<\Bigl(\abs{g'(t_0)}-\tfrac{M}{2}r\Bigr)r,
\end{equation}
则：
\begin{enumerate}
  \item \(f\) 在 \(D(t_0,r)\) 内恰有一个零点（计重数），且该零点为单零点；
  \item 记该零点为 \(\gamma\)，则满足显式位移界
  \begin{equation}\label{eq:xi-rouche-taylor-drift-bound}
  \boxed{\;
  \abs{\gamma-t_0}\ \le\ \frac{\varepsilon}{\abs{g'(t_0)}-\tfrac{M}{2}r}\; }.
  \end{equation}
\end{enumerate}
\end{theorem}
\begin{proof}
令线性化函数 \(\ell(t):=g'(t_0)(t-t_0)\)。
由积分形式的 Taylor 余项，对任意 \(\abs{t-t_0}\le r\) 有
\[
g(t)-\ell(t)=\int_{0}^{1}(1-u)\,g''\!\bigl(t_0+u(t-t_0)\bigr)\,(t-t_0)^{2}\,\dd u,
\]
从而 \(\abs{g(t)-\ell(t)}\le \frac{M}{2}\abs{t-t_0}^{2}\)。
特别地在边界 \(\abs{t-t_0}=r\) 上，
\[
\abs{g(t)-\ell(t)}\le \frac{M}{2}r^{2}<\abs{g'(t_0)}r=\abs{\ell(t)}
\]
由 \eqref{eq:xi-rouche-taylor-conditions} 的 \(\abs{g'(t_0)}>\tfrac{M}{2}r\) 成立。
故由 Rouch\'e 定理，\(g\) 与 \(\ell\) 在 \(D(t_0,r)\) 内的零点数相同，因而 \(g\) 在盘内恰有一个零点（单零点）。
\par
另一方面，对任意 \(\abs{t-t_0}\le r\) 有下界
\[
\abs{g(t)}
\ge
\abs{\ell(t)}-\abs{g(t)-\ell(t)}
\ge
\Bigl(\abs{g'(t_0)}-\tfrac{M}{2}r\Bigr)\abs{t-t_0}.
\]
故在边界 \(\abs{t-t_0}=r\) 上，
\(
\abs{g(t)}\ge\bigl(\abs{g'(t_0)}-\tfrac{M}{2}r\bigr)r>\varepsilon\ge \abs{f(t)-g(t)}
\)
由 \eqref{eq:xi-rouche-taylor-conditions} 的末条不等式成立。
再次应用 Rouch\'e 定理，\(f\) 与 \(g\) 在 \(D(t_0,r)\) 内零点数相同，故 \(f\) 在盘内亦恰有一个零点（计重数），且必须为单零点。
\par
设 \(\gamma\) 为该零点，则 \(f(\gamma)=0\) 蕴含 \(g(\gamma)=-(f-g)(\gamma)\)，从而 \(\abs{g(\gamma)}\le \varepsilon\)。
将前述下界取 \(t=\gamma\) 得
\[
\varepsilon\ \ge\ \abs{g(\gamma)}\ \ge\ \Bigl(\abs{g'(t_0)}-\tfrac{M}{2}r\Bigr)\abs{\gamma-t_0},
\]
即得 \eqref{eq:xi-rouche-taylor-drift-bound}。
\end{proof}

\paragraph{Riemann--Siegel 阶梯递归的量化吸引率}
\begin{corollary}[每加一阶，零点定位误差在主尺度上乘以 \texorpdfstring{$t^{-1/2}$}{t^{-1/2}}]\label{cor:xi-riemann-siegel-attraction-rate}
沿用 \eqref{eq:xi-hardy-Z-def}--\eqref{eq:xi-riemann-siegel-remainder-expansion} 的记号。
设 \(t_0\in\RR\) 且 \(Z^{(m)}(t_0)=0\)、\((Z^{(m)})'(t_0)\neq 0\)。
取 \(r>0\) 使得 \(Z\) 与 \(Z^{(m)}\) 在 \(\overline{D(t_0,r)}\) 上全纯，并记
\[
\varepsilon_m:=\sup_{\abs{t-t_0}\le r}\abs{Z(t)-Z^{(m)}(t)},
\qquad
M_m:=\sup_{\abs{t-t_0}\le r}\abs{(Z^{(m)})''(t)}.
\]
若 \(\varepsilon_m\) 满足余项预算
\begin{equation}\label{eq:xi-riemann-siegel-uniform-remainder-budget}
\varepsilon_m\ \le\ A_m\,a(t_0)^{-m-1/2},
\end{equation}
并且（以 \(f=Z\)、\(g=Z^{(m)}\)）满足定理 \ref{thm:xi-rouche-taylor-zero-transfer} 的条件 \eqref{eq:xi-rouche-taylor-conditions}，则 \(Z\) 在 \(D(t_0,r)\) 内恰有一个单零点 \(\gamma\in\RR\)，并且
\begin{equation}\label{eq:xi-riemann-siegel-zero-drift-budget}
\boxed{\;
\abs{\gamma-t_0}
\ \le\
\frac{A_m\,a(t_0)^{-m-1/2}}{\abs{(Z^{(m)})'(t_0)}-\tfrac{M_m}{2}r}\; }.
\end{equation}
\par
特别地，若进一步具备横截性下界 \(\abs{(Z^{(m)})'(t_0)}\ge B\,t_0^{1/4}\log t_0\)，并取 \(r=r(t_0,m)\) 使得 \(M_m r=o(\abs{(Z^{(m)})'(t_0)})\)，则由 \(a(t_0)^{-m-1/2}\asymp t_0^{-(2m+1)/4}\) 得到量级化吸引率
\begin{equation}\label{eq:xi-riemann-siegel-attraction-scale}
\abs{\gamma-t_0}\ \ll\ \frac{t_0^{-(m+1)/2}}{\log t_0},
\end{equation}
从而递归阶数 \(m\mapsto m+1\) 在相同横截性尺度下使零点定位误差的主量级严格乘上 \(t_0^{-1/2}\)。
\end{corollary}
\begin{proof}[证明要点]
直接以 \(f=Z\)、\(g=Z^{(m)}\) 代入定理 \ref{thm:xi-rouche-taylor-zero-transfer}。
由于 \(Z\) 在实轴取实值，故其零点关于实轴共轭对称；圆盘 \(D(t_0,r)\) 关于实轴对称且仅含一个零点，因而该零点必为实数。
利用 \(a(t_0)^{-m-1/2}=(t_0/(2\pi))^{-(2m+1)/4}\) 与横截性尺度 \(\abs{(Z^{(m)})'(t_0)}\asymp t_0^{1/4}\log t_0\) 化简，得 \eqref{eq:xi-riemann-siegel-attraction-scale}。
\end{proof}

\paragraph{相邻递归层之间零点的同伦对应}
\begin{proposition}[相邻截断层的零点收缩：零点序列在 \texorpdfstring{$m$}{m} 上形成可控 Cauchy]\label{prop:xi-riemann-siegel-adjacent-level-zero-homotopy}
沿用上述记号，并取 \(r>0\) 使得 \(Z^{(m)}\) 与 \(Z^{(m+1)}\) 在 \(\overline{D(t_0,r)}\) 上全纯。
设 \(t_0\) 为 \(Z^{(m)}\) 的单零点，并令
\[
\delta_m:=\sup_{\abs{t-t_0}\le r}\abs{Z^{(m+1)}(t)-Z^{(m)}(t)}.
\]
记 \(M_m:=\sup_{\abs{t-t_0}\le r}\abs{(Z^{(m)})''(t)}\)。
若
\[
\abs{(Z^{(m)})'(t_0)}>\tfrac{M_m}{2}r,
\qquad
\delta_m<\Bigl(\abs{(Z^{(m)})'(t_0)}-\tfrac{M_m}{2}r\Bigr)r,
\]
则存在唯一零点 \(t_1\in D(t_0,r)\) 满足 \(Z^{(m+1)}(t_1)=0\)，并且
\begin{equation}\label{eq:xi-riemann-siegel-adjacent-drift}
\boxed{\;
\abs{t_1-t_0}\ \le\ \frac{\delta_m}{\abs{(Z^{(m)})'(t_0)}-\tfrac{M_m}{2}r}\; }.
\end{equation}
当 \(\delta_m\) 由新增的首个被省略项主导时，\(\delta_m=\Theta(a(t_0)^{-m-1/2})\)，从而在横截性尺度 \(\abs{(Z^{(m)})'(t_0)}\asymp t_0^{1/4}\log t_0\) 下得到
\(
\abs{t_1-t_0}\ll t_0^{-(m+1)/2}/\log t_0
\)。
\end{proposition}
\begin{proof}
以 \(f=Z^{(m+1)}\)、\(g=Z^{(m)}\) 直接套用定理 \ref{thm:xi-rouche-taylor-zero-transfer} 即得 \eqref{eq:xi-riemann-siegel-adjacent-drift}；
最后一条由 \eqref{eq:xi-riemann-siegel-remainder-expansion} 的幂阶提升（每加一阶引入一个额外的 \(a^{-1}\sim t^{-1/2}\) 因子）给出。
\end{proof}

\paragraph{局部误差预算下的零点计数闭合}
\begin{theorem}[区间内零点计数的局部封闭性：屏障带与局部圆盘证书]\label{thm:xi-local-zero-count-closure-budget}\leanverified{paper\_xi\_local\_zero\_count\_closure\_budget}
设 \(I=[T_1,T_2]\subset\RR\)，并设 \(f,g\) 在某个包含 \(I\) 的复邻域内全纯，且对一切 \(t\in I\) 都有 \(f(t),g(t)\in\RR\)。
设 \(g\) 在 \(I\) 内恰有 \(k\) 个单零点
\(
T_1<\tau_1<\cdots<\tau_k<T_2
\)。
取常数 \(\eta>0\) 与半径 \(r_1,\dots,r_k>0\) 满足：
\begin{enumerate}
  \item \(\,[\tau_j-r_j,\tau_j+r_j]\subset (T_1,T_2)\) 且这些区间两两不交；
  \item \(\inf_{t\in I\setminus \bigcup_{j=1}^{k}(\tau_j-r_j,\tau_j+r_j)}\abs{g(t)}\ge 2\eta\)；
  \item \(\sup_{t\in I}\abs{f(t)-g(t)}\le \eta\) 且对每个 \(j\) 有 \(\sup_{\abs{t-\tau_j}\le r_j}\abs{f(t)-g(t)}\le \eta\)；
  \item 对每个 \(j\)，以 \(t_0=\tau_j\) 在圆盘 \(\overline{D(\tau_j,r_j)}\) 上满足定理 \ref{thm:xi-rouche-taylor-zero-transfer} 的条件 \eqref{eq:xi-rouche-taylor-conditions}（取 \(f\) 与 \(g\) 为给定函数）。
\end{enumerate}
则 \(f\) 在 \(I\) 内恰有 \(k\) 个实零点 \(\gamma_1<\cdots<\gamma_k\)，并且对每个 \(j\) 都有唯一对应
\[
\gamma_j\in (\tau_j-r_j,\tau_j+r_j),
\qquad
f(\gamma_j)=0,
\]
且不存在任何额外零点落在 \(I\) 的其余部分。
此外，每个 \(\gamma_j\) 均满足定理 \ref{thm:xi-rouche-taylor-zero-transfer} 的位移界 \eqref{eq:xi-rouche-taylor-drift-bound}（以 \(t_0=\tau_j\) 代入）。
\end{theorem}
\begin{proof}
由条件 (ii)--(iii)，若 \(t\in I\setminus \bigcup_j(\tau_j-r_j,\tau_j+r_j)\)，则
\(
\abs{f(t)}\ge \abs{g(t)}-\abs{f(t)-g(t)}\ge 2\eta-\eta=\eta>0
\)，故 \(f\) 在该集合上无零点。
\par
另一方面，对每个 \(j\)，由条件 (iv) 与定理 \ref{thm:xi-rouche-taylor-zero-transfer}（取 \(t_0=\tau_j\)、半径 \(r_j\)）可知：\(f\) 在 \(D(\tau_j,r_j)\) 内恰有一个单零点；由于 \(f\) 在实轴取实值且圆盘关于实轴对称，该零点必为实数，记为 \(\gamma_j\in(\tau_j-r_j,\tau_j+r_j)\)。
区间不交性保证 \(\gamma_j\) 两两不同。
合并上述两部分即得：\(f\) 在 \(I\) 内恰有 \(k\) 个实零点，且无额外零点。
位移界由定理 \ref{thm:xi-rouche-taylor-zero-transfer} 逐点给出。
\end{proof}

\begin{remark}[计数闭合的制度含义]\label{rem:xi-local-zero-count-closure-meaning}
定理 \ref{thm:xi-local-zero-count-closure-budget} 给出一种纯局部的“计数 + 定位”闭合接口：当候选零点邻域的全纯证书（圆盘上的 Rouch\'e--Taylor 预算）与邻域外的屏障带下界同时成立时，区间内零点的存在性、唯一性与位移界可在同一误差预算链内闭合，而无需诉诸全局型零点计数流程。
\end{remark}

\paragraph{递归校正零点的定向律（猜想）}
\begin{conjecture}[首省略项锁定零点漂移方向]\label{conj:xi-riemann-siegel-omitted-term-direction-law}
设 \(\gamma_n\) 为 Hardy 函数 \(Z(t)\) 的第 \(n\) 个正零点（按大小排列），并令 \(\tau_n^{(m)}\) 为 \(Z^{(m)}\) 的第 \(n\) 个正零点（在一致编号可实现的范围内）。
将 Riemann--Siegel 余项的第 \(m\) 阶修正项记为
\[
\Delta_m(t):=(-1)^{N+1}a(t)^{-m-1/2}C_m\!\bigl(z(t)\bigr),
\]
其中 \(a,N,z\) 由 \eqref{eq:xi-riemann-siegel-aNz} 给出。
则对每个固定 \(n\) 存在 \(m_0(n)\) 使得对所有 \(m\ge m_0(n)\) 成立严格定向律
\[
\operatorname{sgn}\bigl(\tau_n^{(m)}-\gamma_n\bigr)
=
-\,\operatorname{sgn}\bigl(\Delta_m(\tau_n^{(m)})\bigr),
\]
并且序列 \(\{\tau_n^{(m)}\}_{m\ge m_0(n)}\) 对 \(\gamma_n\) 呈现单调收敛或二步交错收敛（由 \(\Delta_m\) 的符号稳定性决定）。
若该猜想成立，则零点递归不仅具有 \(t^{-1/2}\) 量级的几何收缩率（推论 \ref{cor:xi-riemann-siegel-attraction-rate}），且其漂移方向由单一可计算修正项在局部上锁定。
\end{conjecture}

\endinput
